
\documentclass[12pt]{article} % for submission to journals
%\documentclass[11pt]{article}

%%% PAGE DIMENSIONS1
\usepackage{geometry} 
%\geometry{top=1.45in, bottom=1.45in, left=1.25in, right=1.25in} % for submission to journals
%\geometry{top=1.4in, bottom=1.4in, left=1.25in, right=1.25in}
\geometry{top=1.0in, bottom=1.0in, left=1.0in, right=1.0in}

%%% FONTS

%\usepackage{newpxtext}
%\usepackage{newpxmath}

%\usepackage{XCharter}

%\usepackage{palatino}

%\usepackage{tgpagella}

\usepackage[sc]{mathpazo}

\usepackage[T1]{fontenc}

%%% FIGURES
\usepackage{graphicx} % support the \includegraphics command and options
\usepackage{floatrow}
\usepackage{epstopdf}
\usepackage{subfig}
\usepackage{float}
\usepackage{morefloats}
\usepackage[skip=0.5em,font=normalsize,labelfont = bf]{caption} % for submission to journals

%\usepackage[skip=0.5em,font=small]{caption}

%%% PACKAGES
\usepackage{url}
\usepackage[colorlinks=true,
            linkcolor=blue,
            urlcolor=blue,
            citecolor=blue]{hyperref}
\usepackage{rotating}
\usepackage{multirow}
%\usepackage[parfill]{parskip} % Activate to begin paragraphs with an empty line rather than an indent
\usepackage[section]{placeins}
\usepackage{booktabs} % for much better looking tables
\usepackage{tabularx}
\usepackage{lscape}
\usepackage{paralist} % very flexible & customisable lists (eg. enumerate/itemize, etc.)
\usepackage{verbatim} % adds environment for commenting out blocks of text & for better verbatim
\usepackage{setspace}
\usepackage{csquotes}
\usepackage{derivative}
\usepackage{dirtytalk}
\usepackage{graphicx}
\usepackage{adjustbox}

\floatsetup[table]{capposition=top}
\floatsetup[figure]{capposition=top}

\usepackage{tikz}
\usetikzlibrary{arrows,positioning} 
\tikzset{
    %Define standard arrow tip
    >=stealth',
    %Define style for boxes
    punkt/.style={
           rectangle,
           rounded corners,
           draw=black, very thick,
           text width=6.5em,
           minimum height=2em,
           text centered},
    % Define arrow style
    pil/.style={
           ->,
           thick,
           shorten <=2pt,
           shorten >=2pt,}
}


\usepackage[round]{natbib}
\renewcommand\bibname{References}

%\usepackage{footbib}

%\usepackage[style=verbose-ibid,backend=bibtex]{biblatex}
%\bibliography{../JMP_bib.bib}



%%% MATH
\usepackage{amsmath}
\usepackage{amssymb}
\usepackage{amsfonts}
\usepackage{amsthm}
\usepackage{hyperref}
%\usepackage[amsmath]{ntheorem}
\usepackage{array} % for better arrays (eg matrices) in maths
\usepackage{bbm}
\usepackage{mathrsfs}


\usepackage{accents}
\newcommand{\ubar}[1]{\underaccent{\bar}{#1}} % generate lower bar

\def\thickhline{\noalign{\hrule height.8pt}}

% Strikethrough text
\usepackage[normalem]{ulem}
\usepackage[flushleft]{threeparttable}

%%% HEADERS & FOOTERS
%\usepackage{fancyhdr} % This should be set AFTER setting up the page geometry
%\pagestyle{fancy} % options: empty , plain , fancy
%\renewcommand{\headrulewidth}{0pt} % customise the layout...
%\lhead{\emph{Diversification and Boom-Bust Cycles}}\chead{}\rhead{}
%\lfoot{}\cfoot{\thepage}\rfoot{}

% bibliography
\bibliographystyle{plainnat}
\setcitestyle{authoryear,round,semicolon,aysep={}}


%%% SECTION TITLE APPEARANCE
\usepackage{sectsty}
%\allsectionsfont{\sffamily\mdseries\upshape} % (See the fntguide.pdf for font help)
% (This matches ConTeXt defaults)

%%% ToC (table of contents) APPEARANCE
\usepackage[nottoc,notlof,notlot]{tocbibind} % Put the bibliography in the ToC
\usepackage[titles,subfigure]{tocloft} % Alter the style of the Table of Contents

\clubpenalty = 1000
\widowpenalty = 1000

\theoremstyle{theorem}
\newtheorem*{claim}{Claim}
\newtheorem{theorem}{Theorem}
\newtheorem{proposition}{Proposition}
\newtheorem{lemma}{Lemma}
\newtheorem{assumption}{Assumption}
\newtheorem{corollary}{Corollary}
\newtheorem{remark}{Remark}

\theoremstyle{definition}
\newtheorem{example}{Example}
\newtheorem{mydef}{Definition}
\DeclareMathOperator*{\argmin}{arg\,min}
\usepackage{bm}


%\numberwithin{theorem}{section}
%\numberwithin{prop}{section}
%\numberwithin{lemma}{section}
%\numberwithin{remark}{section}
%\numberwithin{corollary}{section}
%\theoremstyle{definition}
%\theoremstyle{remark}

\def\var{{\rm var}}
\def\cov{{\rm cov}}
\def\corr{{\rm corr}}
\def\proj{{\rm proj}}
\def\E{\mathbb{E}}
\def\e{\varepsilon}
\def\d{\partial}
\def\convp{\overset{P}{\to}}
\def\convd{\overset{d}{\to}}
\def\convas{\overset{a.s.}{\to}}
\def\1{\mathbbm{1}}
\def\ci{\perp\!\!\!\perp}
\usepackage{enumitem}
\usepackage{indentfirst}


\begin{document}



\title{Notes on the Code: Demographic Transition and Fiscal Sustainability}

\author{Rodolfo Rigato \& Rafael Lincoln}

\maketitle

\section*{Class HouseholdOLG() - Household Block}

The class \textbf{HouseholdOLG()} contains all the functions relates to the household block of the model. 

\subsection*{Instances}
The instances of the class are:

\begin{itemize}
\item \textbf{a\_grid}: The asset grid. Is an exogenous parameter
\item \textbf{z\_grid}: The productivity grid. Is an exogenous parameter
\item \textbf{Pi}: The transition matrix for idiosyncratic productivity of households. This is determined by the function markov\_chain\_ar1(), and is exogenous
\item \textbf{T \& Tref}: Number of periods lived \& retirement age. Are an exogenous parameter
\item \textbf{inputs}: dictionary of calibrated parameters. All of this calibration is exogenous \\


All of the following variables are determined endogenously
\item \textbf{outputs}: class which returns the outputs solved in the bellman equation
\item \textbf{pol\_interp}: Convex combination of weights to preserve the mass of assets given endogenous grid choice
\item \textbf{EVa}: Expected marginal value of assets to the household
\item \textbf{g}: Age, productivity and wealth distribution
\item \textbf{aggregates}: Dictionary with aggregate variables
\end{itemize}

\subsection*{Functions}

The functions of the class are:

\begin{itemize}
\item \textbf{bellman\_iteration}: This function solves the value function backwards by iterating period after period from the maximum age $T_{ret}$. Start with the $\mathbb{E}[V_a]$ being zero for agents at retirement (they die with certainty), use the FOC to obtain an initial value for consumption

\begin{align} \label{eq:c_foc}
c_{\text{FOC}}(a,z,j) = (u^{\prime})^{-1} \left( \beta (1-p_{j}) \mathbb{E}[V_a(a^{\prime},z^{\prime},j+1)|a,z,j]\right)
\end{align} 

we then take the $c_{\text{FOC}}$ object and solve for the consumption policy function using the endogenous grid point method. Notice a very important thing: \color{red} The marginal utility of assets $\mathbb{E}[V_a(a^{\prime},z^{\prime},j+1)|a,z,j]$ is decreasing in future assets $a^{\prime}$. Therefore, $c_{\text{FOC}}$ is increasing in the quantity of future assets. \color{black} 

After obtaining the consumption policy function at maximum age, $c(a,z,T_{\text{ret}})$, we solve for the marginal value of assets $V_a (a,z,T_{\text{ret}})$ and use this to iterate backwards and solve for other periods. We obtain the vector of indexes and weights of the linear interpolation

\item \textbf{update\_distribution}: This function calculates the stationary-distribution of assets-income-cohorts $g(a,z,j)$ from the initial distribution of newborns. By hypothesis, we assume that the newborn distribution $g(a,z,0) = \delta (a-0,z-med(z))$ is dirac at zero assets and median income, and iterate forward from this distribution to obtain the distribution for other cohorts. To do such, we do a loop in which, for each age $j$ and point in the grid $(z_i,a_k)$, we use the policy functions to send mass to assets next age $j$ $(z_i,a^{\prime}_k)$. Afterwards, use the transition matrix to update the idiosyncratic income shock $z$ to $z^{\prime}$, arriving to the next age in the distribution. We do this until the maximum age before certain death.

\textbf{Obs:} \color{red} In the case of using this function for calculating the transition path of MIT shocks\color{black}, we use the $g_{0}$ argument in the function to start from the previous distribution in the initial period. 

\item \textbf{policy\_ss}: Take policy functions from the Bellman Equation
\item \textbf{distribution\_ss}: Take distribution from the update\_distribution function
\item \textbf{steady\_state}: Set the instances of the class as the steady state objects calculated via the bellman iteration and the update\_distribution
\item \textbf{mit\_shock}: The MIT shock function has the "Mean-Field Game" kind of blocks, with (i) a block for calculating the value function at the new steady state - terminal condition $\mathbb{E}[V_a]$; (ii) a block for iterating backwards the value functions and solving for the policy functions in time; and (iii) a block for iterating the distribution from the initial steady state forward by using the policy function solved in block (ii). However, this function only calculates this path of distribution/aggregates/policy functions.
\end{itemize}

\section*{Other Functions}

These functions are out of the household block, and connect that block with other blocks of the model - constituting the general equilibrium.

\begin{itemize}
\item \textbf{solve\_steady\_state}: This function is responsible for computing the steady state, connecting the household block with the rest of the model. Define the error function:
\begin{align*}
F(z) = \left[\begin{array}{c}
Y - Y(z) \\ 
r - r(z) \\ 
w - w(z) \\ 
G + (r-g^c)B - T(z) \\ 
(1 + g^c)Transfer - \Psi_t
\end{array} \right]
\end{align*}
where $z = (A,\beta,w,\tau,Transfer)$ is the vector of equilibrium objects of the model. To find $z^*$ which clears all markets $F(z^*) = 0$, we use a solver.

\item \textbf{solve_\transition\_path}: This function is responsible for computing the transition path to aggregate MIT-Shocks. The algorithm used is the newton-method, used as follows:
\begin{enumerate}
\item Compute a 
\end{enumerate}
\end{itemize}

\section*{Heterogeneous Agents Utils}

\begin{itemize}
\item \textbf{solve\_euler\_one\_asset}: We solve for the next period asset policy function $a^{\prime}(a,z,j)$ and consumption policy function $c(a,z,j)$ from the Bellman Equation consumption first order condition.
 From the FOC of the bellman equation above, we take each point on the grid state space $(z_i,a_k)$ and iterate with all the possible choices of assets next period $a^{\prime}_{l}$ to find implied consumption
\begin{align} \label{eq:c_here}
c_{here} = \text{budget}(z_i,a_k) - a^{\prime}_{l}, \vspace{2mm} \forall l \in \{1,\ldots,N_{a}\}
\end{align} 
By using equations \eqref{eq:c_foc} and \eqref{eq:c_here} together with the Intermediate Value Theorem, there must exist an $a^*$ such that these equations are alike.
\begin{align*}
c_{here}(a) - c_{\text{FOC}}(a^*) = 0
\end{align*}
take the value in the grid $a_{iap}$ which is closest to this $a^*$ (that is, the first value where $c_{\text{FOC}}(a_{iap}) > c_{here}(a_{iap}) $ ) and we make a linear interpolation between $a_{iap-1}$ and $a_{iap}$ to solve for the asset policy function $a^{\prime}(a,z,j)$. This function also returns the index of the asset policy function, based on the point in the grid $a_{ind}$
\item \textbf{update\_policy\_1a}: Function which takes the previous distribution $g_{t}(a,z)$ for a given cohort of age $j$ to a mid-object which is the distribution next period conditional on no changes in the productivity $g_{mid,t}(a,z)$. It uses the asset policy function and interpolation weights to distribute the mass of $g_{t}(a,z)$ to $g_{mid,t}(a,z)$. 
\item \textbf{update\_exogenous\_1a}: This function takes the object $g_{mid,t}(a,z)$ and uses the rouwenhorst discretization matrix to distribute idiosyncratic income shocks, obtaining the next period distribution $g_{t+1}(a,z)$
\end{itemize}



\end{document}